\documentclass[twocolumn]{iopjournal}

\usepackage{amsmath,amssymb,amsfonts}
\usepackage{amsthm}

\fancyhfoffset[L]{0mm}
\setlength{\marginparwidth}{0pt}
\setlength{\marginparsep}{0pt}

% ==========================================
% ZENODO PREPRINT HEADERS (ACTIVE)
% ==========================================
\renewcommand{\articletype}[1]{%
  {\vspace*{-8mm}\noindent \Large \sf Preprint}\par
  \vspace*{6mm}\noindent{\scriptsize \sf{\bfseries \MakeUppercase{#1}}}}

\fancyhead[L]{}
\fancyhead[R]{}

\hypersetup{pdftitle={Vortices of the Stellated Octahedron}, pdfauthor={Stephen Nelson}, pdfsubject={Gravity Field Dynamics (GFD)}, pdfkeywords={stellated octahedron, vortex, standing wave, tetrahedral symmetry, Field Origin, proof}}

\newtheorem{lemma}{Lemma}
\newtheorem{theorem}{Theorem}
\newtheorem{proposition}{Proposition}
\newtheorem{corollary}{Corollary}
\theoremstyle{definition}
\newtheorem{definition}{Definition}
\theoremstyle{remark}
\newtheorem*{remark}{Remark}

\begin{document}

\articletype{Proof}

\title{The Walsh Eigenmodes of the Pixel, and What Closure Reads}

\author{Stephen Nelson}

\affil{Field Dynamics Research}

\ead{stephen@fielddynamics.org}

\begin{abstract}
The eight Walsh functions $1, x, y, z, yz, xz, xy, xyz$ on the corners of
the cube are known to be its eigenmodes. We ask which of them the GFD
pixel holds, and which of them closure can read. Under the pixel's
symmetry $T_d$ the corners carry $2A_1 \oplus 2T_2$: the \emph{uniform and
parity modes} $1$ and $xyz$ span $A_1$, and the six \emph{direction modes}
span two copies of $T_2$. We prove that every Walsh function is an
eigenmode of a $T_d$-invariant operator exactly when TetA and TetB are
alike (full cube symmetry); otherwise the grades pair, $1 \leftrightarrow
xyz$ and $x \leftrightarrow yz$, $y \leftrightarrow xz$, $z \leftrightarrow
xy$, with mixing angles in closed form. The pixel's vortices lie entirely
in the direction modes, and on the plain cube two degenerate Walsh
standing waves combine into a circulating one. Closure reads the two
sectors separately and completely: every $T_d$-equivariant linear
direction readout of a tetrahedron is Letter~1's frame readout, which by
closure is blind to the uniform mode and reads exactly the direction
modes; every equivariant scalar readout reads only $1$ and $xyz$. In the
archive's oscillator node, the closure's $\mu$ is the TetA share, a
function of $1$ and $xyz$ alone, and its amplitude closure sees only the
total Walsh power, by Parseval. For an eigenmode of the $A_1$ block,
$\mu = x/(1+x)$ holds exactly when the mixing angle satisfies
$\tan\theta = 1/x$. Every statement is checked exactly, or on the node's
own code.
\end{abstract}

\keywords{Walsh functions, Hadamard, eigenmodes, standing wave,
tetrahedral symmetry, closure, stellated octahedron, proof}

% ======================================================================

\section*{Setting}

\textbf{The cube.} The corners are $s \in \{\pm1\}^3$; TetA is the four
with $xyz = +1$ and TetB the four with $xyz = -1$~\cite{NelsonHadamard}.
The Walsh functions are $w_S(s) = \prod_{d\in S}s_d$ for
$S \subseteq \{x, y, z\}$, of \emph{grade} $\lvert S\rvert$; they are
orthogonal, $w_S\cdot w_T = 8\delta_{ST}$. We call $1$ and $xyz$ the
\emph{uniform and parity modes}, and $x, y, z, yz, xz, xy$ the
\emph{direction modes}. On TetA, $yz = x$, $xz = y$, $xy = z$, and the
Walsh matrix factorises as $H_2 \otimes \hat H$, with $H_2$ spatial
parity and $\hat H$ the bordered frame~\cite{NelsonHadamard}.

\textbf{Symmetries.} $F \cong \mathbb Z_2^3$ is the group of coordinate
sign flips, the cube's bit flips. $T_d$ (24 elements) is the group of
signed permutations that keep TetA, and hence TetB; $O_h = T_d \times
\{\pm I\}$, where the inversion $-I$ swaps TetA and TetB. A group acts on
corner amplitudes by $(g\psi)(gs) = \psi(s)$.

\textbf{Operators.} $H$ is a real symmetric operator on the eight corner
amplitudes: a Hamiltonian, or the coupling matrix of eight oscillators,
whose real eigenmodes are standing waves~\cite{NelsonVortex}.

% ======================================================================

\section{The Walsh functions are the cube's eigenmodes}

\begin{theorem}[Known]
\label{thm:walsh}
$H$ commutes with every bit flip if and only if every Walsh function is
an eigenmode of $H$; such $H$ have 8 free parameters, one eigenvalue per
Walsh function. In particular the cube graph has $Aw_S = (3 -
2\lvert S\rvert)w_S$, and its Laplacian $L = 3I - A$ has $Lw_S =
2\lvert S\rvert\,w_S$~\cite{ODonnell2014}. Each $w_S$ is real, so as a
mode of coupled oscillators it is a standing wave and carries no current.
\end{theorem}

\begin{proof}
The $w_S$ are the eight distinct characters of $F$, so $F$ acts on the
corner space without multiplicity, and by Schur's
lemma~\cite{Serre1977} an operator commuting with $F$ is diagonal in the
character basis; conversely a Walsh-diagonal operator commutes with $F$.
A bit flip at $s$ changes $w_S$ by the factor $-1$ for each of the
$\lvert S\rvert$ coordinates in $S$, which gives $A$.
\end{proof}

% ======================================================================

\section{The pixel pairs the grades}

\begin{theorem}[Walsh modes of the pixel]
\label{thm:pairing}
\begin{enumerate}
\item[(i)] Under $T_d$ the corners carry $2A_1 \oplus 2T_2$: $1$ and
  $xyz$ are each invariant and span $A_1$; grade 1 ($x, y, z$) and grade 2
  ($yz, xz, xy$) are each a copy of $T_2$. Multiplying by $xyz$ commutes
  with $T_d$ and maps grade 1 onto grade 2.
\item[(ii)] The most general $T_d$-invariant $H$ has six parameters:
  corner energies $\epsilon_A, \epsilon_B$; face-diagonal couplings
  $h_{AA}, h_{BB}$ inside each chamber; cube-edge and body-diagonal
  couplings $h_e, h_d$ between the chambers. In the chamber basis
  ($u_{A,B}$, the normalised indicators; $f^d_{A,B}$, the coordinate $s_d$
  restricted to a chamber and normalised),
  \begin{equation}
  H = \begin{pmatrix}\alpha_A & \beta\\ \beta & \alpha_B\end{pmatrix}
  \oplus
  \begin{pmatrix}\gamma_A & \delta\\ \delta & \gamma_B\end{pmatrix}^{\oplus 3},
  \label{eq:blocks}
  \end{equation}
  with $\alpha_{A,B} = \epsilon_{A,B} + 3h_{AA,BB}$, $\beta = 3h_e + h_d$,
  $\gamma_{A,B} = \epsilon_{A,B} - h_{AA,BB}$ and $\delta = h_e - h_d$.
  In this basis $1 \propto u_A + u_B$, $xyz \propto u_A - u_B$,
  $x \propto f^x_A + f^x_B$ and $yz \propto f^x_A - f^x_B$.
\item[(iii)] Every Walsh function is an eigenmode of $H$ if and only if
  $\epsilon_A = \epsilon_B$ and $h_{AA} = h_{BB}$, that is, if and only
  if $H$ is $O_h$-invariant (TetA and TetB alike). The eigenvalues are
  then $\alpha \pm \beta$ on $1$ and $xyz$, and $\gamma + \delta$,
  $\gamma - \delta$ on grades 1 and 2.
\item[(iv)] Otherwise $H$ mixes exactly the pairs $1 \leftrightarrow
  xyz$ and $x \leftrightarrow yz$, $y \leftrightarrow xz$, $z
  \leftrightarrow xy$, the pairs related by multiplying by $xyz$. The
  eigenmodes are $(\cos\theta, \sin\theta)$ in the chamber basis of each
  block, with
  \begin{equation}
  \tan 2\theta_{A_1} = \frac{2\beta}{\alpha_A - \alpha_B},\quad
  \tan 2\theta_{T_2} = \frac{2\delta}{\gamma_A - \gamma_B},
  \label{eq:mixing}
  \end{equation}
  and a mode's weight on the lower Walsh grade of its pair is
  $(\cos\theta + \sin\theta)^2/2$.
\end{enumerate}
\end{theorem}

\begin{proof}
(i) By characters: $1$ and $xyz$ are $T_d$-invariant (every element of
$T_d$ keeps the chambers); grade 1 transforms as the coordinates, which is
$T_2$ of $T_d$; $xyz$ is invariant, so multiplication by it is
equivariant and carries grade 1 to grade 2. (ii) $T_d$ has two orbits of
corners and four of corner pairs (distances $2\sqrt2$ inside each
chamber, $2$ and $2\sqrt3$ between them). By Schur, $H$ preserves the
$A_1$ and $T_2$ parts and acts on $T_2$ as a $2\times2$ matrix tensored
with $I_3$; the entries follow from row sums (each corner has three
chamber partners, three edge partners and one diagonal partner) and, on
$T_2$, from $\sum_{b \sim a} s_d(b) = s_d(a)$ along edges and $s_d(-a) =
-s_d(a)$. (iii) In each block, $(1, \pm1)$ are eigenvectors if and only
if the diagonal entries agree, i.e.\ $\alpha_A = \alpha_B$ and $\gamma_A =
\gamma_B$; this is $\epsilon_A = \epsilon_B$, $h_{AA} = h_{BB}$, which is
invariance under $-I$, since $-I$ swaps the chambers and fixes the
edge and diagonal pairs. Under $O_h$ the four grades are four
inequivalent representations ($-I$ gives grade $\lvert S\rvert$ the sign
$(-1)^{\lvert S\rvert}$). (iv) is the eigenproblem of a real symmetric
$2\times2$ matrix.
\end{proof}

\begin{corollary}[Vortices are made of direction modes]
In the 15-site pixel (Field Origin, octahedron, corners)~\cite{NelsonVortex}
the sites carry $4A_1 \oplus E \oplus 3T_2$: two $T_2$ copies on the
corners and one on the octahedron. The rotoreflection $S$ has eigenvalue
$+i$ only on $T_2$, so the winding space $V_{+i}$ takes one vector from
each $T_2$ copy. Every vortex of the pixel is therefore built from the
direction modes: its corner part has no component on $1$ or $xyz$, and
the uniform and parity modes never wind.
\end{corollary}

\begin{proof}
Characters, as in Theorem~\ref{thm:pairing}(i). $S$ acts as the identity
on $A_1$, has eigenvalues $\pm1$ on $E$ ($S^2$ is a rotation by $\pi$,
which acts trivially on $E$), and has eigenvalues $i, -i, -1$ on each
$T_2$.
\end{proof}

\begin{proposition}[Two standing waves make a vortex]
\label{prop:cube}
On the cube graph, $H = tA$ with $t > 0$ (full cube symmetry), $\psi = x +
iy$ is an eigenmode at the doubly degenerate grade-1 energy $t$. It winds
under the quarter turn $C_4$ about $\hat z$, $C_4\psi = -i\psi$, and
carries the circulating current $-8t$ (that is, $-t/2$ per unit norm)
through every half-plane bounded by the $z$ axis, counted as the physical
flow $J_{i\to j} = 2\,\mathrm{Im}(\bar\psi_jH_{ji}\psi_i)$ in the
right-handed sense~\cite{NelsonVortex}. The Walsh modes $x$ and $y$
alone carry none.
\end{proposition}

\begin{proof}
$Ax = x$ and $Ay = y$ (Theorem~\ref{thm:walsh}); $C_4$ sends $x \mapsto y$
and $y \mapsto -x$ on functions, so $C_4\psi = y - ix = -i\psi$. The
current through three half-planes is computed exactly; a real mode has
$\mathrm{Im}(\psi_jH_{ji}\psi_i) = 0$.
\end{proof}

When TetA and TetB differ, the plain quarter turn is no longer a
symmetry and grade 1 mixes with grade 2 (Theorem~\ref{thm:pairing}); the
pixel's vortices then wind under $S$ instead~\cite{NelsonVortex}.

% ======================================================================

\section{Closure reads exactly the direction modes}

\begin{theorem}[What closure reads]
\label{thm:closure}
\begin{enumerate}
\item[(i)] Every linear map $R$ from TetA's four corner amplitudes to
  $\mathbb R^3$ with $R(g\psi) = gR(\psi)$ for all $g \in T_d$ is a
  multiple of Letter~1's frame readout $\psi \mapsto \sum_i \psi_iv_i$,
  where $v_1, \dots, v_4$ are the legs (the TetA corners)~\cite{NelsonL1}.
\item[(ii)] The frame readout annihilates the uniform mode exactly
  because of closure, $v_1 + v_2 + v_3 + v_4 = 0$, and is injective on
  the sum-zero modes. These are the direction modes on TetA, where its
  three rows are $x = yz$, $y = xz$ and $z = xy$.
\item[(iii)] On all eight corners the equivariant direction readouts form
  a two-dimensional space, spanned by the grade-1 readout
  $\sum_k\psi_k s_k$ and the grade-2 readout $\sum_k\psi_k\,(xyz)(s_k)\,
  s_k$, equivalently by the TetA and TetB frames. Both vanish on $1$ and
  $xyz$. Every equivariant scalar readout, conversely, is constant on each
  chamber, so it reads only $1$ and $xyz$.
\end{enumerate}
So the symmetric linear readouts split the eight Walsh functions
completely: directions read exactly the six direction modes, scalars read
exactly the uniform and parity modes, and no symmetric linear readout
does anything else.
\end{theorem}

\begin{proof}
TetA's corners carry $A_1 \oplus T_2$ and $\mathbb R^3$ carries $T_2$, so
$\mathrm{Hom}_{T_d}$ is one-dimensional (Schur), and the frame readout is
equivariant and nonzero. Its kernel is the $A_1$ part, the uniform mode,
and $\sum_i v_i = 0$ is the statement that the uniform mode maps to zero.
On eight corners $\mathrm{Hom}_{T_d}(2A_1 \oplus 2T_2, T_2)$ has
dimension 2 and $\mathrm{Hom}_{T_d}(2A_1 \oplus 2T_2, A_1)$ has dimension
2; the stated readouts are equivariant and independent. All three
statements are also solved directly from the linear equivariance
equations.
\end{proof}

% ======================================================================

\section{What the oscillator node's closure reads}

The archive's octahedral resonance node~\cite{NelsonPixelated} drives
eight corner oscillators, mixes them, and forms a vertex state $v_k$ with
Walsh coefficients $c_S = \frac18\sum_k w_S(s_k)v_k$. Its closure has two
parts: it computes $\mu = (1 + c_{xyz}/c_1)/2$, and it relaxes every
corner amplitude toward a target $a^*$ at rate $\kappa = 0.02$, where $a^*
= r/(1 + r)$ in the bounded rule and $\min(r, 5)$ in the original, with
$r^2$ the mean energy of the two mixes.

\begin{theorem}[The node's closure]
\label{thm:node}
\begin{enumerate}
\item[(i)] $\mu = S_A/(S_A + S_B)$, the TetA share of the vertex state
  ($S_{A,B}$ the sums over each chamber). It depends only on the uniform
  and parity modes, so adding any direction modes, in particular any
  vortex (Corollary~1), leaves it unchanged.
\item[(ii)] The energy that drives the closure is
  $\sum_kv_k^2 = 8\sum_Sc_S^2$ (Parseval), the total Walsh power. Every
  corner has the same target, so while no clamp is active the spread
  between corner amplitudes contracts by exactly $1 - \kappa$ per step,
  whatever the signal; and at equal amplitudes a step multiplies the
  carrier's whole Walsh spectrum by one factor.
\item[(iii)] For an eigenmode $(\cos\theta, \sin\theta)$ of the $A_1$
  block (\ref{eq:blocks}), $\mu = 1/(1 + \tan\theta)$. Hence $\mu =
  x/(1+x)$ exactly when $\tan\theta_{A_1} = 1/x$.
\end{enumerate}
\end{theorem}

\begin{proof}
(i) $c_1 = (S_A + S_B)/8$ and $c_{xyz} = (S_A - S_B)/8$, since $xyz = \pm1$
on the two chambers. (ii) Parseval for the orthogonal Walsh basis. The
update $a_k \leftarrow a_k + \kappa(a^* - a_k)$ with one $a^*$ for all
$k$ gives $a_k - a_j \leftarrow (1-\kappa)(a_k - a_j)$; the carrier is
linear in a common amplitude. (iii) The mode has $S_A = 2\cos\theta$ and
$S_B = 2\sin\theta$. Checked on the node's own code for 800 steps in both
modes.
\end{proof}

% ======================================================================

\section{Remarks}

\begin{remark}[What has been proved]
Within the stated class, three things are now theorems. The pixel holds
the Walsh modes exactly when TetA and TetB are alike, and holds them in
pairs, rotated by (\ref{eq:mixing}), when they differ. Its vortices are
made of the direction modes. And any symmetric linear readout must read
the Walsh spectrum in two sectors: closure of the frame reads the six
direction modes and is blind to the uniform and parity modes, while
$\mu$ reads those two and is blind to the directions. The node's closure
sees the rest of the spectrum only through its total power.
\end{remark}

\begin{remark}[Not claimed]
That the physical field is this node, or closes this way; that its
readouts are linear; or that the field's $A_1$ mixing angle satisfies
$\tan\theta = a_0/g$. Theorem~\ref{thm:node}(iii) turns that last
question into a test: it is what $\mu = x/(1+x)$ requires of the uniform
and parity block.
\end{remark}

% ======================================================================

\section*{Verification}

\texttt{verify\_walsh\_eigenmodes.py} runs 34 checks: exact
representation theory over $T_d$ and $O_h$, the general bit-flip and
$T_d$-invariant operators solved symbolically in all parameters, the
blocks (\ref{eq:blocks}) and the condition of
Theorem~\ref{thm:pairing}(iii) solved exactly, the pixel's
decomposition and winding space, Proposition~\ref{prop:cube} through
three half-planes, the equivariant readouts solved from their linear
equations, and Theorem~\ref{thm:node} both symbolically and on the node's
code (\texttt{octahedral\_node.py} in~\cite{NelsonPixelated}).

\begin{thebibliography}{9}

\bibitem{NelsonHadamard}
S.~Nelson,
\textit{The Stellated Octahedron, the Cube and the Hadamard Matrix},
Field Dynamics Research (2026),
\url{proofs/hadamard-cube-frame}.

\bibitem{NelsonVortex}
S.~Nelson,
\textit{Vortices of the Stellated Octahedron},
Field Dynamics Research (2026),
\url{proofs/vortices-of-the-stellated-octahedron}.

\bibitem{NelsonL1}
S.~Nelson,
\textit{The Coordinate-Space Geometry of the Tetrad Frame} (Letter~1),
Field Dynamics Research (2026).

\bibitem{NelsonPixelated}
S.~Nelson,
\textit{The Universe Is Pixelated}, research note with the octahedral
resonance node, Field Dynamics Research (2026),
\url{research/universe-is-pixelated}.

\bibitem{ODonnell2014}
R.~O'Donnell,
\textit{Analysis of Boolean Functions}
(Cambridge University Press, Cambridge, 2014).

\bibitem{Serre1977}
J.-P.~Serre,
\textit{Linear Representations of Finite Groups}
(Springer, New York, 1977).

\end{thebibliography}

\end{document}
